% GENERATED FILE. Edit canonical lessons, then run npm run generate:books.
% canonical-source-hash: fnv1a64:c29157b335bcfb71
% canonical-lessons: ZH-C88-liang, ZH-W88-liang, ZH-C88-yi-sandhi, ZH-C88-bei, ZH-C88-ben, ZH-C88-zhang, ZH-C88-zhi, ZH-C88-practice

\chapter{Two Cups of Tea}
\label{ch:zh-two-cups-of-tea}
\hlchaptermodality{88}

\begin{chapteropening}
\textbf{By the end of this chapter:} \emph{I can count things with \zh{两} and the right measure word.}

Counting things with \zh{两} and four measure words beyond \zh{个}, and the tone \zh{一} takes before each, built on one new written character: \zh{两}.
\end{chapteropening}

\section[\texorpdfstring{\zh{两} (liǎng)}{两 (liǎng)}]{\zh{两} (liǎng) — two of something}

\label{lesson:ZH-C88-liang}



\begin{quote}
A new chapter, and it opens with a second word for two.

\begin{itemize}
\raggedright
  \item \emph{Recall:} say \emph{nǐ yǒu méiyǒu yàoshi}, and answer \emph{méiyǒu} — \textbf{R1}, the chapter you have finished
  \item \emph{Recall:} say \zh{一个人} — \emph{yí ge rén}, one person
  \item \emph{Recall:} say \zh{二} — \emph{èr}, two, as you count
\end{itemize}
\end{quote}

\subsection*{You'll want to know: \zh{两}}
\textbf{\zh{两}} — \emph{liǎng} — \textquotedblleft{}two (of something)\textquotedblright{}.

\begin{quote}
\textbf{\zh{两个人}} \emph{liǎng ge rén} — two people
\end{quote}

The next lesson puts the new character in your hand.

\begin{cousinweb}
\textbf{\zh{两}} (\emph{liǎng}) is inherited from Middle Chinese \emph{ljangX} (Guangyun, in Baxter's transcription), where it names a weight and, quoting an older dictionary, 'twice'. 'Two of a thing' grew from those senses by a step the Guangyun does not show.
\end{cousinweb}

\begin{grammarlens}[title={What you've built}]
Counting aloud, two is \zh{二}: \emph{yī, èr, sān}. In front of a measure word it is \zh{两}: \zh{两个人}, two people. \emph{èr ge rén} is not said.
\end{grammarlens}

\subsection*{Your turn}
\begin{itemize}
\raggedright
  \item \emph{Say it:} \emph{liǎng}
  \item \emph{Say it:} \emph{liǎng}, once more
  \item \emph{Say it:} \zh{两个人} — \emph{liǎng ge rén} — two people
  \item \emph{Say it:} \zh{一} \zh{二} \zh{三} — \emph{yī èr sān} — counting: one, two, three
  \item \emph{Recall:} say \emph{yí ge rén}, then \emph{liǎng ge rén}
\end{itemize}

\begin{tcolorbox}[breakable,colback=teal!4,colframe=teal!35!black,title={Before you move on}]
When is two \zh{两}? (\textbf{In front of a measure word: \zh{两个人}}.) And when is it \zh{二}? (\textbf{Counting: \emph{yī, èr, sān}}.)
\end{tcolorbox}
\section[\texorpdfstring{\zh{两} (liǎng)}{两 (liǎng)}]{\zh{两} — a bar, a frame, two small \zh{人}}

\label{lesson:ZH-W88-liang}



\begin{quote}
The word first, then the character inside it.

\begin{itemize}
\raggedright
  \item \emph{Recall:} say \emph{liǎng} — \textbf{R1}, one lesson back, before it decays
\end{itemize}

\begin{quote}
\zh{一} \zh{人}
\end{quote}

\emph{Read it:} each one aloud, with nothing said first — \emph{yī}, \emph{rén}
\end{quote}

\subsection*{Script — a bar, a frame, two small \zh{人}}
\hlblockfigure{\detokenize{figures/ZH-W88-liang-filmstrip.pdf}}{How \zh{两} is written, stroke by stroke}

\begin{quote}
\zh{两}
\end{quote}

\textbf{A bar on top, then a frame, then what the frame holds.}

\textbf{The top bar, one stroke}, across.

\textbf{The frame, two strokes}: a short vertical on the left, then a horizontal that turns down the right side and hooks to the left.

\textbf{Two small \zh{人}, four strokes}, the left one first: each a left-falling stroke and a short stroke falling to the right, rising through the top of the frame.

1 + 2 + 4 = seven.

\subsection*{Your turn}
\begin{itemize}
\raggedright
  \item \emph{Write it:} \zh{两}, all seven, with the model in front of you
  \item \emph{Cover:} write it again from memory, and count the pieces, not the strokes
\end{itemize}

\begin{tcolorbox}[breakable,colback=teal!4,colframe=teal!35!black,title={Before you move on}]
How many strokes in \zh{两}? (\textbf{Seven}.) What does the frame hold? (\textbf{Two small \zh{人}}.) Which comes first, the frame or the two \zh{人}? (\textbf{The frame}.)

Source: \href{https://raw.githubusercontent.com/chanind/hanzi-writer-data/68d10a4b21150cae5e1ebbd223eed289cf32d90c/data/\%E4\%B8\%A4.json}{Hanzi Writer Data \zh{两}, pinned stroke paths and medians}.
\end{tcolorbox}
\section[yī to yí ge, yì tiān (yí ge, yì tiān)]{yī $\to$ yí ge, yì tiān: the tone \zh{一} takes}

\label{lesson:ZH-C88-yi-sandhi}



\begin{quote}
Before the measure words, the sound that changes in front of them.

\begin{itemize}
\raggedright
  \item \emph{Read it:} \zh{两个人}, aloud — \emph{liǎng ge rén}
  \item \emph{Recall:} say \zh{一百} — \emph{yìbǎi}, a hundred
  \item \emph{Recall:} say \zh{一起} — \emph{yìqǐ}, together
\end{itemize}
\end{quote}

\begin{sounds}
\zh{一} has three spoken tones, and the next syllable picks one.

\textbf{Counting, and inside numbers and dates, it stays level}: \emph{yī, èr, sān}; \textbf{\zh{星期一}} \emph{xīngqīyī}.

\textbf{Before a fourth tone it rises}: \textbf{\zh{一个}} \emph{yí ge}. \zh{个} is fourth tone, even when it is said light.

\textbf{Before the other three tones it falls}: \textbf{\zh{一天}} \emph{yì tiān}, \textbf{\zh{一年}} \emph{yì nián}, \textbf{\zh{一百}} \emph{yìbǎi}.

This book writes the tone you say, as it writes \emph{bú shì} for \zh{不是}.
\end{sounds}

\subsection*{Your turn}
Say these aloud:

\begin{itemize}
\raggedright
  \item \emph{yī, èr, sān} — level
  \item \zh{一个} — \emph{yí ge} — rising, before a fourth tone
  \item \zh{一天} — \emph{yì tiān} — falling, before a first tone
  \item \zh{一年} — \emph{yì nián} — falling, before a second tone
  \item \zh{一百} — \emph{yìbǎi} — falling, before a third tone
\end{itemize}

\begin{tcolorbox}[breakable,colback=teal!4,colframe=teal!35!black,title={Before you move on}]
What does \zh{一} do before a fourth tone? (\textbf{It rises: \emph{yí ge}}.) Before the other tones? (\textbf{It falls: \emph{yì tiān}}.) And counting? (\textbf{It stays level: \emph{yī}}.)
\end{tcolorbox}
\section[bēi]{bēi — a cup of, a glass of}

\label{lesson:ZH-C88-bei}



\begin{quote}
Before the new one: the tone of \zh{一}, then two words that go with this one.

\begin{itemize}
\raggedright
  \item \emph{Recall:} say \zh{一个} — \emph{yí ge}, then \zh{一天} — \emph{yì tiān}
  \item \emph{Recall:} say \emph{chá}, tea — \textbf{R4}, recognition at distance
  \item \emph{Recall:} say \emph{bēizi}, a cup — \textbf{R3}, durable
\end{itemize}
\end{quote}

\subsection*{You'll want to know: bēi}
\emph{bēi} — \textquotedblleft{}a cup of; a glass of\textquotedblright{}, a measure word.

\begin{quote}
\emph{liǎng bēi chá} — two cups of tea
\end{quote}

\begin{quote}
\emph{yì bēi shuǐ} — a glass of water
\end{quote}

Its character comes later; for now it is a word for the ear.

\begin{cousinweb}
\emph{bēi} is inherited from Middle Chinese \emph{pwoj} (Guangyun, in Baxter's transcription), glossed through an older book as a cup. It is the bēi of \emph{bēizi}; as a measure word it counts what a cup holds.
\end{cousinweb}

\begin{grammarlens}[title={What you've built}]
A measure word counts by the container: \emph{liǎng bēi chá}, two cups' worth. \zh{一} before its level tone falls: \emph{yì bēi}.
\end{grammarlens}

\subsection*{Your turn}
\begin{itemize}
\raggedright
  \item \emph{Say it:} \emph{bēi}
  \item \emph{Say it:} \emph{bēi}, once more
  \item \emph{Say it:} \emph{liǎng bēi chá} — two cups of tea
  \item \emph{Say it:} \emph{yì bēi shuǐ} — a glass of water
  \item \emph{Recall:} say \emph{bēizi}, then \emph{bēi} again
\end{itemize}

\begin{tcolorbox}[breakable,colback=teal!4,colframe=teal!35!black,title={Before you move on}]
What does \emph{bēi} count? (\textbf{Cups and glasses of something}.) Ask for two cups of tea. (\emph{liǎng bēi chá}.)
\end{tcolorbox}
\section[běn]{běn — counting books}

\label{lesson:ZH-C88-ben}



\begin{quote}
Before the new one: last lesson's word, then the thing this one counts.

\begin{itemize}
\raggedright
  \item \emph{Recall:} say \emph{liǎng bēi chá}, two cups of tea — \textbf{R1}, one lesson back
  \item \emph{Recall:} say \zh{书} — \emph{shū}, a book
\end{itemize}
\end{quote}

\subsection*{You'll want to know: běn}
\emph{běn} — a measure word for books.

\begin{quote}
\emph{yì běn shū} — a book
\end{quote}

\begin{quote}
\emph{wǒ yǒu liǎng běn shū} — I have two books.
\end{quote}

Its character comes later; for now it is a word for the ear.

\begin{cousinweb}
\emph{běn} is inherited from Middle Chinese \emph{pwonX} (Guangyun, in Baxter's transcription), glossed 'the root; the beginning'. Counting books with it is a later use the Guangyun does not record.
\end{cousinweb}

\begin{grammarlens}[title={What you've built}]
\zh{个} counts almost anything; \emph{běn} is the word a speaker reaches for with books. \zh{一} before its dip falls: \emph{yì běn}.
\end{grammarlens}

\subsection*{Your turn}
\begin{itemize}
\raggedright
  \item \emph{Say it:} \emph{běn}
  \item \emph{Say it:} \emph{běn}, once more
  \item \emph{Say it:} \emph{yì běn shū} — a book
  \item \emph{Say it:} \emph{wǒ yǒu liǎng běn shū} — I have two books.
  \item \emph{Recall:} say \emph{liǎng bēi chá}, then \emph{liǎng běn shū}
\end{itemize}

\begin{tcolorbox}[breakable,colback=teal!4,colframe=teal!35!black,title={Before you move on}]
What does \emph{běn} count? (\textbf{Books}.) Say you have two. (\emph{wǒ yǒu liǎng běn shū}.)
\end{tcolorbox}
\section[zhāng]{zhāng — counting flat things}

\label{lesson:ZH-C88-zhang}



\begin{quote}
Before the new one: last lesson's word, then two flat things.

\begin{itemize}
\raggedright
  \item \emph{Recall:} say \emph{yì běn shū}, a book — \textbf{R1}, one lesson back
  \item \emph{Recall:} say \emph{piào}, a ticket — \textbf{R3}, durable
  \item \emph{Recall:} say \emph{chuáng}, a bed — \textbf{R1}, from the chapter before
\end{itemize}
\end{quote}

\subsection*{You'll want to know: zhāng}
\emph{zhāng} — a measure word for flat things: tables, beds, tickets.

\begin{quote}
\emph{yì zhāng zhuōzi} — a table
\end{quote}

\begin{quote}
\emph{liǎng zhāng piào} — two tickets
\end{quote}

Its character comes later; for now it is a word for the ear.

\begin{cousinweb}
\emph{zhāng} is inherited from Middle Chinese \emph{trjang} (Guangyun, in Baxter's transcription), glossed 'to stretch out; to spread'. Counting things that lie spread flat is a later use, not given there.
\end{cousinweb}

\begin{grammarlens}[title={What you've built}]
A table, a bed, a ticket: things with a flat top or a flat face take \emph{zhāng}.
\end{grammarlens}

\subsection*{Your turn}
\begin{itemize}
\raggedright
  \item \emph{Say it:} \emph{zhāng}
  \item \emph{Say it:} \emph{zhāng}, once more
  \item \emph{Say it:} \emph{liǎng zhāng piào} — two tickets
  \item \emph{Say it:} \emph{yì zhāng chuáng} — a bed
  \item \emph{Recall:} say \emph{yì běn shū}, then \emph{yì zhāng zhuōzi}
\end{itemize}

\begin{tcolorbox}[breakable,colback=teal!4,colframe=teal!35!black,title={Before you move on}]
What does \emph{zhāng} count? (\textbf{Flat things: tables, beds, tickets}.) Ask for two tickets. (\emph{liǎng zhāng piào}.)
\end{tcolorbox}
\section[zhī]{zhī — counting animals}

\label{lesson:ZH-C88-zhi}



\begin{quote}
Before the new one: last lesson's word, then two animals.

\begin{itemize}
\raggedright
  \item \emph{Recall:} say \emph{liǎng zhāng piào}, two tickets — \textbf{R1}, one lesson back
  \item \emph{Recall:} say \emph{māo}, a cat — \textbf{R2}
  \item \emph{Recall:} say \emph{gǒu}, a dog — \textbf{R3}
\end{itemize}
\end{quote}

\subsection*{You'll want to know: zhī}
\emph{zhī} — a measure word for animals, and for one of a pair.

\begin{quote}
\emph{liǎng zhī māo} — two cats
\end{quote}

\begin{quote}
\emph{yì zhī shǒu} — one hand
\end{quote}

Its character comes later; for now it is a word for the ear.

\begin{cousinweb}
\emph{zhī} is inherited from Middle Chinese \emph{tsyek} (Guangyun, in Baxter's transcription), glossed 'one'; the Guangyun, quoting an older dictionary, describes the old character as a hand holding a single bird.
\end{cousinweb}

\begin{grammarlens}[title={What you've built}]
Cats, dogs, birds: animals take \emph{zhī}. So does one of a pair, such as a hand: \emph{yì zhī shǒu}.
\end{grammarlens}

\subsection*{Your turn}
\begin{itemize}
\raggedright
  \item \emph{Say it:} \emph{zhī}
  \item \emph{Say it:} \emph{zhī}, once more
  \item \emph{Say it:} \emph{liǎng zhī māo} — two cats
  \item \emph{Say it:} \emph{wǒ jiā yǒu yì zhī gǒu} — We have a dog at home.
  \item \emph{Recall:} say \emph{liǎng zhāng piào}, then \emph{liǎng zhī māo}
\end{itemize}

\begin{tcolorbox}[breakable,colback=teal!4,colframe=teal!35!black,title={Before you move on}]
What does \emph{zhī} count? (\textbf{Animals, and one of a pair}.) Say two cats. (\emph{liǎng zhī māo}.)
\end{tcolorbox}
\section[Practice]{\zh{两个人}}

\label{lesson:ZH-C88-practice}



\begin{quote}
The chapter closes here, with two people looking for a table.

\begin{itemize}
\raggedright
  \item \emph{Recall:} say \emph{liǎng zhī māo}, two cats — \textbf{R1}, one lesson back
  \item \emph{Recall:} say \emph{shàngbian}, on top — \textbf{R2}
  \item \emph{Recall:} say \emph{jǐ}, how many — \textbf{R4}, recognition at distance
\end{itemize}
\end{quote}

\subsection*{The exchange}
\begin{quote}
\textbf{A:} \emph{nǐmen jǐ ge rén} — How many of you are there?
\end{quote}

\begin{quote}
\textbf{B:} \zh{两个人}
\end{quote}

\begin{quote}
\emph{liǎng ge rén} — Two.
\end{quote}

\begin{quote}
\textbf{A:} \zh{你们喝什么}
\end{quote}

\begin{quote}
\emph{nǐmen hē shénme} — What will you drink?
\end{quote}

\begin{quote}
\textbf{B:} \emph{liǎng bēi chá} — Two cups of tea.
\end{quote}

\begin{quote}
\textbf{A:} \emph{hǎo, nà zhāng zhuōzi} — Fine. That table, then.
\end{quote}

\begin{quote}
\textbf{B:} \emph{nàr yǒu liǎng zhī māo} — There are two cats there.
\end{quote}

\begin{quote}
\textbf{A:} \emph{zhuōzi shàngbian yǒu yì běn shū} — And there's a book on the table.
\end{quote}

Every count has its measure word: \zh{个} for people, \emph{bēi} for tea, \emph{zhāng} for the table, \emph{zhī} for the cats, \emph{běn} for the book. And every two is \zh{两}.

\begin{tcolorbox}[breakable,skin=enhanced,colback=orange!6,colframe=orange!45!black,boxrule=0.5pt,arc=1mm,left=6pt,right=6pt,top=4pt,bottom=4pt,fonttitle=\bfseries,title={Read this}]
\emph{Read it:} the exchange through once without stopping, then answer the question below
\end{tcolorbox}

\begin{grammarlens}[title={What you've built — a number, a measure word, a thing}]
Every count has three parts, in this order: the number, the measure word, the thing.

\begin{itemize}
\raggedright
  \item \textbf{\zh{两个人}} \emph{liǎng ge rén} — two people
  \item \emph{yì bēi shuǐ} — a glass of water
  \item \emph{yì běn shū} — a book
  \item \emph{liǎng zhāng piào} — two tickets
  \item \emph{liǎng zhī māo} — two cats
\end{itemize}

\textbf{Listen.} \emph{Nǐ yǒu jǐ běn shū?}

\textbf{Speak.} Count what is on your table, with the right measure word for each.

\emph{Write it:} \zh{两} — the bar, the frame, then the two small \zh{人}
\end{grammarlens}

\subsection*{Your turn}
\begin{itemize}
\raggedright
  \item \emph{Say it:} \zh{两个人} — \emph{liǎng ge rén} — two people
  \item \emph{Say it:} \emph{liǎng bēi chá} — two cups of tea
  \item \emph{Say it:} \emph{yí ge}, \emph{yì běn}, \emph{yì zhāng}, \emph{yì zhī} — and listen to \zh{一} move
  \item \emph{Say it:} the whole exchange, both parts
  \item \emph{Write it:} \zh{两个人}
\end{itemize}

\begin{tcolorbox}[breakable,colback=teal!4,colframe=teal!35!black,title={Before you move on}]
Two people? (\textbf{\zh{两个人}}.) Two cups of tea? (\emph{liǎng bēi chá}.) A book? (\emph{yì běn shū}.) Two cats? (\emph{liǎng zhī māo}.)
\end{tcolorbox}
